\documentclass[11pt, a4paper]{article}

% ---------- PACKAGES ----------
\usepackage[top=1.2cm, bottom=1.2cm, left=1.5cm, right=1.5cm]{geometry}
\usepackage{titlesec}
\usepackage{enumitem}
\usepackage{hyperref}
\usepackage{xcolor}
\usepackage{fontspec}
\usepackage{microtype}

\usepackage[normalem]{ulem}  % For underline
\setmainfont{Times New Roman}

% ---------- COLORS (for links only) ----------
\hypersetup{
  colorlinks=true,
  urlcolor=blue
}

\pagestyle{empty}
\setlength{\parindent}{0pt}
\setlength{\parskip}{1.5pt}
\setlength{\emergencystretch}{3em}  % Avoid overfull hbox
\definecolor{sectiongray}{RGB}{60,60,60}


% ---------- SECTION STYLE ----------
\definecolor{sectiongray}{RGB}{60,60,60}

\titleformat{\section}
  {\large\bfseries\color{sectiongray}}
  {}
  {0em}
  {}
  [\vspace{2pt}\titlerule]

\titlespacing{\section}{0pt}{10pt}{6pt}



% ---------- LIST STYLE ----------
\setlist[itemize]{
    leftmargin=1.2em,
    itemsep=1.5pt,
    topsep=2pt,
    parsep=0pt
}

% ---------- DOCUMENT ----------
\begin{document}

% ================= HEADER =================
\begin{center}
{\LARGE \textbf{BRYAN STEVEN MENOSCAL SANTANA}}
\end{center}

\vspace{6pt}

\begin{tabular*}{\textwidth}{@{\extracolsep{\fill}} l l}
\textbf{Ubicación:} Ecuador, Manabí, Manta & \textbf{Teléfono:} +593 99 739 9441 \\
\textbf{Correo:} \href{mailto:bryanmenoscal2005@gmail.com}{bryanmenoscal2005@gmail.com}
& \textbf{Portafolio:} \href{https://stevsant.vercel.app?ref=cv-backend-es}{stevsant.vercel.app} \\
\textbf{LinkedIn:} \href{https://linkedin.com/in/bryanmenoscal26}{linkedin.com/in/bryanmenoscal26}
& \textbf{GitHub:} \href{https://github.com/StevSant}{github.com/StevSant} \\
\end{tabular*}

\vspace{6pt}
\hrule
\vspace{8pt}


% ================= PERFIL =================
\section*{PERFIL PROFESIONAL}

Ingeniero Backend especializado en el diseño y desarrollo de sistemas distribuidos y APIs escalables con Python (FastAPI, Django REST). Experiencia construyendo arquitecturas basadas en microservicios y eventos (Kafka), aplicando principios de Clean Architecture y desplegando servicios en entornos cloud. Fuerte énfasis en calidad de software mediante estrategias de testing integral (unitarias, integración y E2E) con PyTest.
% ================= IDIOMAS =================
\section*{IDIOMAS}
\vspace{-2pt}
\begin{tabular}{p{0.25\textwidth} p{0.7\textwidth}}
\textbf{Español:} & Nativo \\
\textbf{Inglés:} & C1
\end{tabular}
\vspace{6pt}
\hrule
\vspace{8pt}

% ================= SKILLS =================
\section*{HABILIDADES TÉCNICAS}
\vspace{-2pt}
\begin{tabular}{p{0.25\textwidth} p{0.7\textwidth}}
\textbf{Backend:} & Django, Django REST Framework, Flask, FastAPI, NestJS, Gin, Kafka, RabbitMQ, LangChain, LangGraph \\
\textbf{Lenguajes:} & Python (principal), Golang (básico), TypeScript \\
\textbf{Cloud \& DevOps:} & AWS (boto/aioboto), Docker, Google Cloud Run, Azure, Vercel \\
\textbf{Observabilidad:} & OpenTelemetry, Prometheus, Grafana \\
\textbf{Frontend:} & Angular, React (conocimientos básicos) \\
\textbf{Bases de datos:} & MySQL, PostgreSQL, MongoDB \\
\textbf{Testing:} & PyTest, Pruebas Unitarias, Pruebas de Integración, Pruebas End-to-End (E2E) \\
\textbf{Herramientas:} & Git, Docker, Kubernetes, GitHub / GitLab, Postman, Visual Studio Code \\
\textbf{Metodologías:} & Scrum, Desarrollo Ágil, DevOps, TDD \\
\end{tabular}
\vspace{6pt}
\hrule
\vspace{8pt}

% ================= EXPERIENCIA =================
\section*{EXPERIENCIA PROFESIONAL}

\textbf{PATRIMORE} \hfill Remoto \\
\textit{AI Agent Developer Intern} \hfill Mar 2026 -- May 2026

\begin{itemize}
\item Desarrollador de agentes autónomos asistidos por IA utilizando LangGraph, FastAPI y los SDKs de Anthropic (incluyendo Model Context Protocol, MCP).
\item Integración de servicios con AWS mediante boto y aioboto.
\item Implementación de herramientas de observabilidad del sistema utilizando OpenTelemetry, Prometheus y Grafana.
\end{itemize}

\vspace{4pt}

% ================= PROYECTOS =================
\section*{PROYECTOS Y EXPERIENCIA RELEVANTE}

\textbf{SME RISK ASSESSMENT AGENT --- HACKATÓN VIAMATICA (2025)} \hfill Ecuador, Guayas, Guayaquil \\
\textit{Backend \& AI Developer | Hackathon (460+ participantes, 3er lugar)} \hfill 2025

\begin{itemize}
\item Desarrollé el agente de IA con LangChain y la API REST con FastAPI para la comunicación con el frontend y la evaluación de riesgo crediticio en tiempo real para PYMES.
\item Implementé pipelines de web scraping automatizado con Playwright para extraer datos financieros de la Superintendencia de Compañías y el SRI, además de datos de reputación digital de Twitter, Facebook, Instagram y LinkedIn.
\item Diseñé un sistema de scoring de reputación digital que incorpora la presencia online de las PYMES en la evaluación crediticia, abordando sesgos en los métodos tradicionales de evaluación.
\item Escribí pruebas unitarias y de integración con PyTest para validar respuestas del agente de IA, pipelines de scraping y endpoints de la API bajo escenarios de borde.
\end{itemize}

\vspace{4pt}

\textbf{STREAMFLOWMUSIC} \hfill Proyecto Personal \\
\textit{Backend Developer} \hfill 2025

\begin{itemize}
\item Desarrollé una API REST escalable utilizando Django REST Framework siguiendo principios de arquitectura limpia.
\item Implementé tareas asíncronas en segundo plano para la ingesta de contenido usando la API de YouTube.
\item Optimicé consultas a la base de datos para soportar más de 100 usuarios concurrentes con baja latencia.
\item Implementé un suite de pruebas integral con PyTest cubriendo pruebas unitarias, de integración y E2E para endpoints de la API, tareas en segundo plano y operaciones de base de datos.
\end{itemize}

\vspace{4pt}

\textbf{MESAYA} \hfill Proyecto Académico \\
\textit{Backend Developer} \hfill 2025 -- 2026

\begin{itemize}
\item Diseñé una plataforma de reservas de restaurantes basada en microservicios con más de 7 servicios usando NestJS, FastAPI, Go y Python, siguiendo Clean Architecture y patrones orientados a eventos con Apache Kafka.
\item Desarrollé autenticación (JWT), procesamiento de pagos con integración de Stripe, chatbot con IA usando LangChain/Groq, API GraphQL y servicio en tiempo real con WebSockets y webhooks B2B seguros (HMAC-SHA256).
\item Diseñé una arquitectura escalable con PostgreSQL y Redis para manejar reservas concurrentes y flujos automatizados con n8n para notificaciones por correo e integraciones con APIs de terceros.
\item Construí suites de pruebas E2E y de integración con PyTest a través de los microservicios, validando comunicación inter-servicios vía Kafka, flujos de pago y manejo de eventos WebSocket.
\end{itemize}

% ================= EDUCACION =================
\section*{EDUCACIÓN}

\textbf{UNIVERSIDAD LAICA ELOY ALFARO DE MANABÍ} \hfill Ecuador, Manabí, Manta \\
\textit{Ingeniería de Software} \hfill 2023 -- Presente

\begin{itemize}
\item Desarrollo de proyectos académicos y personales utilizando Django REST, Flask, FastAPI, NestJS y Angular, enfocados en resolver problemas reales con soluciones escalables.
\item Formación complementaria autodidacta en backend, frontend y metodologías ágiles (Scrum / Kanban), reforzada con experiencia práctica.
\item \textbf{Cursos relevantes:} Estructuras de Datos, Arquitectura de Software, Bases de Datos, Desarrollo Backend, Testing de Software.
\end{itemize}

% ================= LIDERAZGO =================
\section*{LIDERAZGO Y VOLUNTARIADO}

\textbf{Artificial Intelligence Club -- ULEAM} \hfill Ecuador, Manabí, Manta \\
\textit{Content Management Lead} \hfill 2024 -- Presente

\begin{itemize}
\item Planifiqué y estructuré programas de formación en inteligencia artificial para la comunidad universitaria.
\item Apoyé la organización de charlas técnicas y talleres, fomentando el aprendizaje colaborativo.
\end{itemize}

\end{document}
